\chapter{Validation}\label{ch:validation}

A simulator is only useful if its numbers can be trusted. LinkStudio is validated at three levels:
\begin{itemize}
  \item unit tests of the engine and the DSP helpers;
  \item \emph{physics tests}, which compare complete chains with closed-form results;
  \item a test that loads and runs every shipped scenario.
\end{itemize}
The demo scripts repeat the physics checks in readable form, and this manual's figures are a
third, visual check of the same models.

\section{Running the checks}
\begin{lstlisting}[style=shell]
cd backend
python -m pytest                 # all tests (about 220, a few minutes)
python -m pytest tests/test_physics.py -q
python scripts/main.py           # demos: simulated vs theory, PASS/FAIL per check
python scripts/main.py --quick   # fewer errors per point, faster
\end{lstlisting}
Physics tests stop after 300--500 errors and accept a relative deviation of 20--25\,\%. By
\cref{sec:mc} this is several standard deviations, so false alarms are negligible, while a
systematic error of 1\,dB (a factor of 2--3 in BER on steep curves) is caught.

\section{Physics checks}
\begin{center}
\small
\begin{longtable}{@{}p{3.3cm}p{5.6cm}p{5.3cm}@{}}
\toprule
\textbf{Model} & \textbf{Reference} & \textbf{Checked in}\\
\midrule
\endhead
AWGN, PSK/QAM/PAM & exact Gray BER (\cref{sec:ber-awgn}) & \file{test_physics.py}, \cref{fig:ber-awgn}\\
Soft demapping & max-log decisions = hard decisions & \file{test_physics.py}\\
Pulse shaping & RRC link BER = symbol-rate theory; noise metadata after MF & \file{test_physics.py}\\
OFDM over AWGN & per-subcarrier $E_s/N_0$ theory & \file{test_physics.py}\\
OFDM multipath & LS and perfect CSI equalization & \file{test_physics.py}, \cref{fig:ofdm-chest}\\
Rayleigh + ZF & \cref{eq:rayleigh-ber} & \file{test_physics.py}, \cref{fig:diversity}\\
ISI equalizers & ZF/MMSE/LMS remove ISI & \file{test_physics.py}\\
Soft vs hard Viterbi & soft better by $\approx2$\,dB & \file{test_physics.py}\\
Synchronization & receiver close to theory with delay and CFO; fails without sync & \file{test_advanced.py}\\
Carrier recovery & tracks phase noise & \file{test_advanced.py}\\
Timing recovery & estimates the fractional delay & \file{test_advanced.py}\\
LDPC & codewords satisfy $\mat H\vect c\T=0$; corrects noise; coding gain & \file{test_advanced.py}\\
Polar & noiseless decoding; frozen set; coding gain & \file{test_advanced.py}\\
Puncturing & rate 3/4; decodable & \file{test_advanced.py}\\
SIMO MRC & \cref{eq:mrc}, $L=2$ & \file{test_advanced.py}, \cref{fig:diversity}\\
Spatial multiplexing & ZF $2\times2$ = SISO Rayleigh; ML $<$ MMSE $<$ ZF & \file{test_advanced.py}, \cref{fig:smux}\\
Alamouti & MRC $L=2$ shifted by 3\,dB & \file{test_advanced.py}, \cref{fig:diversity}\\
Time-varying multipath & OFDM with Jakes taps & \file{test_advanced.py}\\
Line codes & round trips; 4B/5B and 64B/66B rates; error multiplication & \file{test_wireline_optical.py}\\
Real NRZ over AWGN & BPSK theory & \file{test_wireline_optical.py}\\
Cable & insertion loss equals the preset formula & \file{test_wireline_optical.py}, \cref{fig:cable}\\
SerDes & CTLE peaking; CTLE + DFE open the eye & \file{test_wireline_optical.py}, \cref{fig:ctle}\\
DMT & water-filling bit loading & \file{test_wireline_optical.py}, \cref{fig:dmt}\\
Crosstalk, echo & coupling level; canceller reduces echo & \file{test_wireline_optical.py}\\
Jitter & zero jitter is identity; RJ degrades & \file{test_wireline_optical.py}\\
Fiber loss, metadata & $P\,e^{-\alpha L}$, $D\ell$ & \file{test_wireline_optical.py}\\
Dispersion & Gaussian broadening; exact compensation & \file{test_wireline_optical.py}, \cref{fig:dispersion}\\
SPM & CW phase $\gamma PL_{\mathrm{eff}}$ & \file{test_wireline_optical.py}, \cref{fig:spm}\\
Split-step solver & second-order convergence on the soliton & \cref{fig:soliton}\\
EDFA, multi-span link & OSNR formulas ($n_{sp}=F/2$, $GA\ne1$, course exercise), saturation & \file{test_optical_models.py}, \cref{fig:osnr-spans,fig:osnr-span-length}\\
Raman & on/off gain, lumped-limit ASE, $F_{\mathrm{eq}}<3$\,dB & \file{test_optical_models.py}, \cref{fig:raman}\\
GN model & long-span form, 3\,dB rule, link metrics; split-step comparison & \file{test_optical_models.py}, \cref{fig:gn-launch}\\
Lasers & threshold and slope, settling, linearized response, adiabatic chirp, VCSEL overlaps, thermal rollover, hole-burning turn-off & \file{test_lasers.py}, \cref{fig:vcsel-li,fig:vcsel-transient,fig:laser-response}\\
Modulators, detectors & MZM sine law and ER, PAM levels, TIA output, APD power law, ISNR & \file{test_optical_models.py}, \cref{fig:snr-receivers}\\
PM-QAM & course BER formulas vs exact; back-to-back simulation & \file{test_optical_models.py}, \cref{fig:coherent-ber}\\
IM/DD OOK, PAM-4 & Q-factor BER; PAM-M formula of the course; dispersion limit & \file{test_wireline_optical.py}, \file{test_optical_models.py}, \cref{fig:pam-imdd,fig:cd-limit}\\
IM/DD PAM-4 & dispersion-induced power fading & \file{test_wireline_optical.py}\\
Coherent DP-QPSK & BER vs OSNR theory within $\approx0.3$\,dB; FO estimate & \file{test_wireline_optical.py}\\
FSO & mean and variance of the scintillation & \file{test_wireline_optical.py}, \cref{fig:fso}\\
WDM & XPM lowers the MER of the centre channel & \file{test_wireline_optical.py}\\
Multimode fiber & $\alpha_{\mathrm{opt}}$, $\sigma_{\min}$, WKB delay, continuum limit, class bandwidths, launch closed forms, variance additivity & \file{test_multimode.py}, \cref{fig:mmf-profile,fig:mmf-launch}\\
Data-center lane & Bessel--Thomson noise bandwidth; reach by fiber class & \file{test_multimode.py}, \cref{fig:mmf-eyes}\\
Fading generator & Zheng--Xiao: $R_h=J_0$, uncorrelated I/Q & \file{test_channels.py}, \cref{fig:jakes}\\
NR TDL profiles & unit normalized delay spread; TDL-D $K_1=13.3$\,dB; fractional delay & \file{test_channels.py}\\
Path loss & P.525; TR~38.901 breakpoint continuity, NLOS $\ge$ LOS, shadowing, LOS probability & \file{test_propagation.py}, \cref{fig:pathloss}\\
Thermal noise & $k_BTF f_s$; free-space QPSK link budget BER & \file{test_propagation.py}\\
Copper loops & G.996.1 Mid-C insertion loss within 0.35\,dB; causal dielectric & \file{test_wireline_optical.py}, \cref{fig:pic}\\
Dispersion slope & group delay $D\ell\Delta\lambda+\frac{S\ell}2\Delta\lambda^2$; CD+slope compensation & \file{test_wireline_optical.py}\\
PMD & Maxwellian DGD (mean, $\langle\tau^2\rangle/\langle\tau\rangle^2$) & \file{test_wireline_optical.py}, \cref{fig:pmd}\\
Adaptive split step & soliton error $<10^{-3}$ with default criterion & \file{test_wireline_optical.py}\\
DML, coherent RX & adiabatic chirp frequency step; LO split per polarization & \file{test_wireline_optical.py}\\
FSO channel & $\alpha,\beta$ from $\sigma_R^2$; Kim, geometric and pointing losses & \file{test_wireline_optical.py}\\
CP-based OFDM sync & recovers delay and CFO & \file{test_wireline_optical.py}\\
Confidence interval & Wilson bounds contain the BER & \file{test_wireline_optical.py}\\
All scenarios & load, validate, run; BER below a threshold & \file{test_scenarios.py}\\
\bottomrule
\end{longtable}
\end{center}

\section{Known limitations}
Every model in this manual states its assumptions; the most important simplifications are
collected here.
\begin{itemize}
  \item Laser dynamics are small-signal (no large-signal rate equations). The multimode fiber
  omits mode-partition and modal noise.
  \item FSO fading is quasi-static (no time correlation, aperture averaging or beam wander).
  \item WDM simulates the full multi-channel field but demodulates only one channel.
  \item The OFDM CP synchronizer estimates the fractional CFO only.
  \item DMT adapts bits to a flat transmit PSD (no power water-filling).
  \item MIMO channels are flat; frequency-selective MIMO-OFDM is not yet a block.
  \item Subsystems are edited in their YAML files, not in the GUI.
  \item Error counts assume independent errors; for coded and block-fading links, also judge
  the accuracy by the number of frame errors.
\end{itemize}
